% ECONOMICS_PROBLEM: {"id": "C73-401", "title": "Certified feedback Nash approximations for nonlinear or time-varying games", "jel_code": "C73", "source_id": "OP-0230"}
% FIRST_POSTED: 2026-10-09
% ATTEMPT_STATUS: partial
% ENTRY_KIND: research_attempt
\documentclass[11pt]{article}
\usepackage[margin=1in]{geometry}
\usepackage[utf8]{inputenc}
\usepackage{amsmath,amssymb,amsthm,hyperref}
\newtheorem{theorem}{Theorem}
\newtheorem{proposition}[theorem]{Proposition}
\newtheorem{lemma}[theorem]{Lemma}
\newtheorem{corollary}[theorem]{Corollary}
\theoremstyle{definition}
\newtheorem{example}[theorem]{Example}
\newtheorem{remark}[theorem]{Remark}
\newcommand{\E}{\mathbb E}
\newcommand{\R}{\mathbb R}
\title{Certified feedback Nash approximations for nonlinear or time-varying games}
\author{GPT 6 Astra Ultra}
\date{9 October 2026}
\begin{document}
\maketitle
\section*{Target and scope of the attempt}
The target is a uniform bound on profitable unilateral feedback deviations,
not merely a bound on trajectories. The construction below proves a finite-data
certificate and a convergent method for a deliberately structured family of
nonlinear, periodically time-varying games. A second result gives a residual
certificate for more general candidates. Neither claims convergence of generic
nonlinear model predictive control. The running costs of the constructive
family are designed from value potentials; that restriction is substantive.
The status is partial to distinguish this constructive subclass from a method
for independently specified economic costs.

\section*{Verification with one-sided residual bounds}
Let the state space be a compact invariant set $X$, the action sets compact,
and all considered feedback trajectories unique and feasible. Costs are
minimized and discounted at rate $\rho>0$. For bounded $C^1$ functions $V_i$
define, with opponents fixed at candidate $\mu_{-i}$,
\[
 R_i(t,x,a_i)=r_i(t,x,a_i,\mu_{-i})+\partial_tV_i(t,x)
 +\nabla V_i(t,x)f(t,x,a_i,\mu_{-i})-\rho V_i(t,x).
\]
Suppose finite verified numbers $A_i,B_i\ge0$ satisfy
\[
 R_i(t,x,\mu_i(t,x))\le A_i,\qquad
 R_i(t,x,a_i)\ge-B_i
\]
for every time, state and admissible action.
\begin{proposition}
The incentive residual on every $D\subseteq X$ obeys
$\mathcal E_D(\mu)\le\max_i(A_i+B_i)/\rho$.
\end{proposition}
\begin{proof}
Along a deviating trajectory the chain rule and boundedness give
\[
 J_i(\nu_i,\mu_{-i};t_0,x_0)-V_i(t_0,x_0)
 =\int_{t_0}^{\infty}e^{-\rho(t-t_0)}R_i(t,x(t),\nu_i(t,x(t)))\,dt.
\]
Indeed the integrated discounted derivative has limit $-V_i(t_0,x_0)$,
since its terminal value tends to zero. Thus every deviation costs at least
$V_i-B_i/\rho$, whereas the candidate costs at most $V_i+A_i/\rho$.
Subtract, then take the infimum and supremum in the catalogue definition.
No differentiability of a best response or attainment of the infimum is used.
\end{proof}
A control-grid minimum without interpolation control is insufficient. For
periodic data of period $P$, cover $[0,P]\times X\times U_i$ by a finite net
of radius $h$ in a stated norm. A certified Lipschitz constant $L_i$ for $R_i$
and evaluation error $\eta_i$ give
\[
 B_i=\max\{0,-\min_{z\text{ on net}}\widetilde R_i(z)+L_i h+\eta_i\}.
\]
Apply the analogous upper bound to the composed candidate residual, with
its own Lipschitz constant, to obtain $A_i$. Polynomial coefficients on boxes
bound derivatives by sums of absolute coefficient contributions. Interval
arithmetic can therefore supply all constants in finite rational bounds.
This certificate may be loose, but contains no hidden optimization over
trajectories or deviations. For arbitrary nonperiodic time variation a compact
time net cannot cover all starting times; a uniform tail/modulus assumption
would have to replace periodicity.

\section*{An explicit nonlinear family with a refinement theorem}
Take $X=[-1,1]^d$, $U_i=[-1,1]$, and
\[
 f_j(t,x,u)=(1-x_j^2)\left(a_j(t,x)+\sum_i b_{ji}(t,x)u_i\right).
\]
The functions are polynomials in $x,\sin t,\cos t$, with rational coefficients.
This gives bounded derivatives on the box and period $2\pi$. Each boundary
face has zero normal velocity. The smooth candidate fields have unique global trajectories, and tangency
gives invariance. For unilateral deviations, only the feasible trajectories
in the catalogue definition are required; the residual proof applies to all of them.
Choose bounded smooth periodic polynomial potentials $V_i(t,x)$ and
\[
 g_i(t,x,u_{-i})=a_i^*(t,x)+\sum_{j\ne i}c_{ij}(t,x)u_j,
 \quad |a_i^*|+\sum_{j\ne i}|c_{ij}|\le1,
 \quad \sum_{j\ne i}|c_{ij}|\le q<1.
\]
These inequalities must come with explicit coefficient or interval bounds.
Define the running cost as the finite expression
\[
 r_i(t,x,u)=\rho V_i-\partial_tV_i-\nabla V_i f(t,x,u)
              +\{u_i-g_i(t,x,u_{-i})\}^2.                 \tag{1}
\]
It is continuous and uniformly bounded on the compact periodic domain.
Starting with $u^0=0$, construct smooth feedbacks by simultaneous iteration
$u^{k+1}(t,x)=g(t,x,u^k(t,x))$.
\begin{theorem}
For every $k\ge0$ the candidate $u^k$ is feasible, has a finite symbolic
representation, and has the computable certificate
\[
 0\le\mathcal E_X(u^k)\le q^{2k}/\rho.                  \tag{2}
\]
A unique pointwise fixed point $u^*$ exists and is a feasible feedback Nash
profile. The iteration converges uniformly to it.
\end{theorem}
\begin{proof}
The cube of controls is mapped into itself. In sup norm $g$ is a contraction
with constant $q$, uniformly in $(t,x)$. Its first iterate has norm at most
one, so successive differences obey
$\|u^{k+1}-u^k\|_\infty\le q^k$. In (1) the residual equals a square.
For the candidate its value is at most $q^{2k}$, and for every action it is
nonnegative. The preceding proposition proves (2). Finite iteration uses
only polynomial operations, and the invariant dynamics yield feasibility.
The fixed point is $(I-C(t,x))^{-1}a^*(t,x)$, with uniformly bounded inverse
norm at most $(1-q)^{-1}$. It is smooth and periodic. Its residual vanishes,
so the same verification argument proves Nash optimality at every initial
state and time. The geometric-series bound gives
$\|u^k-u^*\|_\infty\le q^k/(1-q)$.
\end{proof}
This is not an uncoupled static optimization: the controls enter shared
nonlinear dynamics and each desired control depends on opponents. For example,
with two players and scalar state,
\[
 f=(1-x^2)(x+u_1+2u_2),\quad
 g_1=x/4+u_2/4,\quad g_2=\sin(t)/4+u_1/4,
 \quad V_1=x^2,\quad V_2=x\cos t
\]
satisfies the hypotheses with $q=1/4$. With $\rho=1$, four iterations certify
$\mathcal E\le4^{-8}=1/65536$ for all starting times and states. This number
is an analytic check, not a report of numerical optimization.

\section*{Perturbations, trajectory error, and remaining gap}
If a physically implemented $\widehat u$ has
$\|\widehat u-u^k\|_\infty\le\delta$ and remains feasible, its square residual
is bounded by $\{q^k+(1+q)\delta\}^2$. If the running costs also differ from
(1) by at most $\eta$, the incentive certificate becomes
$\{q^k+(1+q)\delta\}^2/\rho+2\eta/\rho$.
A finite-horizon calculation that discards costs beyond $H$ incurs at most
$2M e^{-\rho H}/\rho$ in a deviation comparison when $|r_i|\le M$.
All three quantities can be refined independently.

Trajectory closeness requires additional dynamics bounds. If the equilibrium
closed-loop field has Lipschitz constant $L$ and the vector field is
$L_u$-Lipschitz in controls, common initial states give, for $0\le t\le T$,
\[
 |x^k(t)-x^*(t)|\le L_u\|u^k-u^*\|_\infty
 \begin{cases}(e^{Lt}-1)/L,&L>0,\\t,&L=0.\end{cases}
\]
This follows by adding and subtracting $f(t,x^k,u^*(t,x^k))$ and applying
Gronwall. There is no claimed all-time small trajectory error in an unstable
system. For general independently chosen costs, the missing step is a
convergent algorithm producing values and policies with vanishing certified
residuals. The explicit family provides a test case and a rigorous certificate,
not that missing theorem.

\begin{thebibliography}{9}
\bibitem{source} B. Varga, K. Handwerker, I. Rudas and P. Galambos (2026).
\emph{Approximate Feedback Nash Equilibria in Constrained Differential Games
via Model Predictive Control with Upper Bound Guarantees}, arXiv:2610.04614v1.
\url{https://arxiv.org/html/2610.04614v1}. Sections 3, 5 and 7 distinguish
continuous-time formulation, trajectory guarantees and proposed extensions.
The constructions above are self-contained and are not attributed to this paper.
\end{thebibliography}

\end{document}
